\documentclass[11pt]{article}
\usepackage[margin=1in]{geometry}
\usepackage{amsmath, amssymb}
\usepackage{enumitem}
\usepackage{titlesec}
\usepackage{hyperref}
\usepackage{booktabs}
\usepackage{xcolor}
\usepackage{parskip}

\hypersetup{colorlinks=true, linkcolor=blue, urlcolor=blue}

\title{\textbf{EE627 --- Class Summary}\\[0.4em]
       \large February 6, 2026}
\author{Prof.\ Rensheng Wang}
\date{}

\begin{document}
\maketitle
\tableofcontents
\newpage

%------------------------------------------------------------
\section{AI Usage Policy Clarification}
%------------------------------------------------------------

A student Q\&A at the start of class clarified the course's stance on AI tools.
The key points are worth recording as they reflect the professor's broader philosophy.

\begin{itemize}
    \item \textbf{AI use is encouraged} --- it is the industry trend and will be
      universal in future careers.
    \item \textbf{The line between acceptable and unacceptable:} copy, then
      \emph{think}. Read the AI output, evaluate it, tune parameters, catch
      mistakes, and note what you learned. Simple copy-paste with no reflection
      is what should be avoided.
    \item \textbf{How grading works:} assignments require sequential, multi-step
      tasks. Even with AI, most students cannot finish them without understanding
      what they are asking for. The sequential guidance itself is the demonstration
      of capability.
    \item \textbf{The career analogy:} In the near future you will be a
      \emph{manager of AI agents}. You do not write the code --- you give clear
      commands and judge whether the output is correct and useful. That judgment
      is what this class trains.
    \item \textbf{Internet analogy from lecture:} When the professor first arrived
      in the US, he went to a travel agent in person to buy a plane ticket --- even
      though the internet already existed. The lesson: having a powerful tool does
      not automatically change how you think. You have to consciously learn to
      leverage it.
\end{itemize}


%------------------------------------------------------------
\section{Correlation: Motivation and Intuition}
%------------------------------------------------------------

\subsection{Why Correlation?}
Correlation is the \textbf{starting point} for nearly all time series analysis, and
a fundamental diagnostic tool in industry data work. The professor's framing:

\begin{itemize}
    \item Given any dataset, your boss asks: \textit{``Help me identify why this
      is happening.''}
    \item Step one is always: compute correlations between candidate explanatory
      variables and the outcome.
    \item \textbf{High positive correlation} $\Rightarrow$ likely root cause
      (e.g., CPU utilization correlates with server downtime).
    \item \textbf{High negative correlation} $\Rightarrow$ likely mitigation lever.
    \item About 80\% of root-cause identification can be done this way before
      reaching for more complex models.
\end{itemize}

\subsection{Why You Remove the Mean First}
The correlation formula subtracts the mean from both variables before forming the
cross-product. This is not a technicality --- it has a direct interpretation.

\begin{itemize}
    \item The mean is a fixed baseline. What we care about is
      \emph{fluctuation around the mean} --- the dynamic, stochastic movement.
    \item A large mean says nothing about predictability. A tiny mean with huge
      variance is actually very informative.
    \item After mean-removal, the expected value of each term is zero. The
      cross-product then purely measures whether $X$ and $Y$ \emph{move together}:
      when $(X - \mu_X) > 0$ and $(Y - \mu_Y) > 0$, or both negative, you get
      a positive product --- coordinated movement.
    \item If they are uncorrelated: positives and negatives cancel on average
      $\Rightarrow$ cross-product averages to zero.
\end{itemize}

\subsection{Why Normalize by the Standard Deviations?}
Dividing by $\sigma_X \sigma_Y$ produces a \textbf{dimensionless} number between
$-1$ and $+1$.

Without normalization, the raw covariance is hard to interpret: is a covariance of
200 large or small? It depends entirely on the scale of the variables.
Normalization gives a universal flag:

\begin{center}
\begin{tabular}{@{}cl@{}}
\toprule
\textbf{Value of $r$} & \textbf{Interpretation} \\
\midrule
$r \approx +1$        & Strong positive: $X$ up $\Rightarrow$ $Y$ up \\
$r \approx -1$        & Strong negative: $X$ up $\Rightarrow$ $Y$ down \\
$r \approx 0$         & No linear relationship; knowing $X$ tells you nothing about $Y$ \\
\bottomrule
\end{tabular}
\end{center}

\subsection{Correlation Is Not Causation}
\textit{Classic example from lecture:} In summer, both ice cream sales and reported
crime events are high. They are correlated. But buying ice cream does not cause
crime. The hidden driver is \emph{temperature}: it causes both, independently.

The correct interpretation of correlation:
\begin{itemize}
    \item Correlation gives you a \textbf{candidate pool} of causes.
    \item Dependence (causation) implies high correlation, but high correlation
      does not imply dependence.
    \item Use correlation to identify \emph{possible} root causes --- then
      investigate the mechanism.
\end{itemize}

\textit{Other examples from lecture:} smoking and cancer (high correlation used
to guide medical advice); real estate prices rising (correlated with other
investment opportunities); war breaking out (correlated with oil, grocery, and
bond prices).


%------------------------------------------------------------
\section{Stationarity: The Prerequisite for Time Series Modeling}
%------------------------------------------------------------

\subsection{Why Stationarity Matters}
Time series analysis is ultimately about \textbf{forecasting}: learning patterns
from the past to predict the future. For that to be possible, the statistical
properties of the series must be stable over time --- otherwise what you learned
from the past is irrelevant to the future.

\subsection{Strict vs.\ Weak Stationarity}
\textbf{Strict stationarity:} all statistical moments (mean, variance, skewness,
kurtosis, \ldots) are time-invariant. Formally: the joint distribution of
$\{X_{t_1}, \ldots, X_{t_k}\}$ is identical to that of any time-shifted window.

\textbf{Weak (2nd-order) stationarity} --- what is used in practice:
\begin{itemize}
    \item Mean $E(X_t) = \mu$ is constant.
    \item Covariance $\text{Cov}(X_t, X_{t-l}) = v_l$ depends only on lag $l$,
      not on $t$.
\end{itemize}

\textit{Intuition from lecture:} Penn Station foot traffic, measured week by week,
shows nearly identical patterns (Monday rush, weekend dip). That regularity is
stationarity. Stock prices, by contrast, can trend, crash, or spike --- clearly
not stationary in their raw form.

\subsection{What to Do When Data Is Not Stationary}
Raw financial time series almost never satisfy stationarity. The answer is
\textbf{preprocessing before modeling}:
\begin{itemize}
    \item Take \emph{returns} (differences or log-differences) instead of raw
      prices --- this removes the trending mean.
    \item Use rolling standardization to stabilize variance.
\end{itemize}
This is why the class starts with stationarity: it is the necessary condition
before any of the models that follow can be applied.


%------------------------------------------------------------
\section{Auto-Correlation}
%------------------------------------------------------------

\subsection{From Correlation to Auto-Correlation}
Standard correlation measures the relationship between \emph{two different}
variables $X$ and $Y$.

\textbf{Auto-correlation} asks the same question about \emph{one variable with
itself, shifted in time}: how correlated is today's value $X_t$ with the value
$l$ days ago, $X_{t-l}$?

Formally, setting $X \to X_t$ and $Y \to X_{t-l}$ in the correlation formula:
\[
r_l = \frac{\text{Cov}(X_t,\, X_{t-l})}{\text{Var}(X_t)}
    = \frac{v_l}{v_0}, \qquad
v_l = E[(X_t - \mu)(X_{t-l} - \mu)]
\]
Properties: $r_0 = 1$,\; $|r_l| \leq 1$,\; $r_l = r_{-l}$.

\subsection{What Auto-Correlation Means Practically}

\textit{From the lecture:} If you know the stock price $l$ days ago, $r_l$ tells
you what fraction of that deviation from the mean carries forward to today.

The best linear predictor of $X_t$ given $X_{t-l}$ is:
\[
X_t - \mu \;\approx\; r_l\,(X_{t-l} - \mu)
\]

\begin{itemize}
    \item If $r_l = 0$: yesterday tells you nothing; your best forecast is simply
      the mean $\mu$.
    \item If $r_l > 0$: yesterday above average $\Rightarrow$ today probably above
      average.
    \item If $r_l < 0$: yesterday above average $\Rightarrow$ today probably below
      average (mean-reverting behavior).
\end{itemize}

\subsection{Recursive Forecasting}
Auto-correlation enables a chain of predictions:
\begin{enumerate}
    \item Know yesterday $X_{t-1}$ $\Rightarrow$ predict today $\hat{X}_t$.
    \item Use $\hat{X}_t$ $\Rightarrow$ predict tomorrow $\hat{X}_{t+1}$.
    \item Repeat for any horizon.
\end{enumerate}

This recursive structure is the foundation for financial derivatives pricing,
bond valuation, and options pricing --- all of which require forecasting future
values from current ones.


%------------------------------------------------------------
\section{The Auto-Regressive (AR) Model}
%------------------------------------------------------------

\subsection{Definition}
An \textbf{AR($p$)} model states that today's value is a linear combination of
the past $p$ days' values, plus random noise:
\[
X_t = a_0 + a_1 X_{t-1} + a_2 X_{t-2} + \cdots + a_p X_{t-p} + \varepsilon_t
\]

\begin{itemize}
    \item \textbf{``Auto'':} the predictor and the predicted variable are the
      \emph{same} series (Apple stock predicting Apple stock).
    \item \textbf{``Regressive'':} history is used to predict the present/future.
    \item $\varepsilon_t$: zero-mean random noise capturing whatever is unexplained
      by history.
\end{itemize}

\subsection{AR(1): Properties}
AR(1) model: $X_t = a_0 + a_1 X_{t-1} + \varepsilon_t$.

\subsubsection{Mean}
Taking expectations of both sides, using stationarity ($E(X_t) = E(X_{t-1}) = \mu$):
\[
\mu = a_0 + a_1 \mu \quad\Rightarrow\quad \mu = \frac{a_0}{1 - a_1}
\]
\textbf{Requirement:} $a_1 \neq 1$. If $a_1 = 1$ (random walk: today = yesterday
+ noise), the mean is infinite and the series is \emph{non-stationary}.

\subsubsection{Variance}
\[
\text{Var}(X_t) = \frac{\sigma_\varepsilon^2}{1 - a_1^2}
\]
\textbf{Requirement:} $a_1^2 < 1$, i.e., $|a_1| < 1$, so variance is
non-negative and finite. This is the \textbf{stationarity condition} for AR(1).

\subsubsection{ACF Decay (Key Result)}
The autocorrelation of an AR(1) series decays \textbf{exponentially}:
\[
r_l = a_1^l
\]

\textit{Intuition from lecture:} If $a_1 = 0.5$, today and yesterday are 50\%
correlated. Today and 10 days ago: $(0.5)^{10} \approx 0.001$ --- essentially
no correlation. The further back you look, the less history matters.

\subsubsection{Stationarity Condition Summary}
\[
|a_1| < 1
\]
If $a_1 = 1$: random walk --- variance grows without bound.\\
If $|a_1| > 1$: explosive --- variance grows to infinity.\\
Only $|a_1| < 1$ gives a well-behaved, mean-reverting, stationary process.

\subsection{AR(1) ACF Interpretation}

\begin{center}
\begin{tabular}{@{}lll@{}}
\toprule
\textbf{$a_1$} & \textbf{ACF behavior} & \textbf{Series character} \\
\midrule
$a_1 = 0.9$    & Slow exponential decay & Very persistent (yesterday predicts today well) \\
$a_1 = 0.5$    & Moderate decay          & Semi-persistent \\
$a_1 = 0.1$    & Fast decay to zero      & Nearly memoryless \\
$a_1 = -0.5$   & Alternating decay       & Mean-reverting (oscillates around mean) \\
\bottomrule
\end{tabular}
\end{center}


%------------------------------------------------------------
\section{Selecting the AR Order $p$: PACF}
%------------------------------------------------------------

\subsection{The Problem}
In practice the order $p$ of an AR model is unknown. You cannot just guess.
Two principled approaches exist:
\begin{enumerate}
    \item \textbf{Partial Autocorrelation Function (PACF)} --- the primary tool.
    \item \textbf{Information criteria} (AIC, BIC) --- covered later.
\end{enumerate}

\subsection{What Is the PACF?}
The PACF answers: \textit{``What is the direct contribution of $X_{t-l}$ to $X_t$,
after removing the influence of all intermediate lags?''}

Consider fitting AR models of increasing order:
\begin{align*}
X_t &= a_{0,1} + a_{1,1}X_{t-1} + \varepsilon_{1t}                                      & \text{AR(1)}\\
X_t &= a_{0,2} + a_{1,2}X_{t-1} + a_{2,2}X_{t-2} + \varepsilon_{2t}                     & \text{AR(2)}\\
X_t &= a_{0,3} + a_{1,3}X_{t-1} + a_{2,3}X_{t-2} + a_{3,3}X_{t-3} + \varepsilon_{3t}   & \text{AR(3)}\\
&\vdots
\end{align*}

The \textbf{lag-$l$ sample PACF} is the coefficient $\hat{a}_{l,l}$ from the
AR($l$) regression --- the coefficient on the \emph{last added lag} in each model.

\begin{itemize}
    \item Lag-1 PACF $= \hat{a}_{1,1}$: direct effect of $X_{t-1}$ on $X_t$.
    \item Lag-2 PACF $= \hat{a}_{2,2}$: \emph{added} contribution of $X_{t-2}$
      beyond what AR(1) already explains.
    \item Lag-3 PACF $= \hat{a}_{3,3}$: added contribution of $X_{t-3}$ beyond
      AR(2), and so on.
\end{itemize}

Each model is estimated by ordinary least-squares (multiple linear regression).

\subsection{Using PACF to Determine $p$}
Key property for a stationary Gaussian AR($p$) model:
\begin{itemize}
    \item $\hat{a}_{p,p}$ converges to the true $a_p$ as sample size $T
      \to \infty$.
    \item $\hat{a}_{l,l} \to 0$ for all $l > p$ --- lags beyond $p$ add nothing.
    \item Asymptotic variance of $\hat{a}_{l,l}$ is $1/T$ for $l > p$.
\end{itemize}

\textbf{Decision rule:} The standard error of the sample PACF is $\approx
\sqrt{1/T}$. Any lag $l$ whose PACF satisfies $|\hat{a}_{l,l}| > 2/\sqrt{T}$
(95\% significance) is a statistically meaningful contributor. The last such lag
is your estimate of $p$.

\subsection{Worked Example}
A time series of length $T = 996$. Standard error $= \sqrt{1/996} \approx 0.032$.
The 95\% significance threshold is $\approx 2 \times 0.032 = 0.064$.

\begin{center}
\begin{tabular}{@{}lrrrrrrrrrrrr@{}}
\toprule
\textbf{Lag}  & 1 & 2 & 3 & 4 & 5 & 6 & 7 & 8 & 9 & 10 & 11 & 12 \\
\midrule
\textbf{PACF} & 0.115 & 0.030 & 0.102 & 0.033 & 0.062 & $-$0.050
              & 0.031 & 0.052 & 0.063 & 0.005 & 0.005 & 0.011 \\
\bottomrule
\end{tabular}
\end{center}

Lags exceeding the threshold ($>0.064$): lags 1, 3, and 9.
\begin{itemize}
    \item At the \textbf{5\% level}: AR(3) or AR(9) are both candidates.
    \item At the \textbf{1\% level} (threshold $\approx 0.081$): only lag 1 clearly
      stands out $\Rightarrow$ specify \textbf{AR(3)}.
\end{itemize}

In practice, prefer the simpler model unless there is a strong theoretical reason
to include the higher-order lag.


%------------------------------------------------------------
\section{AR($p$) Stationarity: Characteristic Equation}
%------------------------------------------------------------

\subsection{The Characteristic Equation}
The stationarity of an AR($p$) model is determined by the roots of its
\textbf{characteristic equation}. Starting from:
\[
X_t = a_1 X_{t-1} + a_2 X_{t-2} + \cdots + a_p X_{t-p} + \varepsilon_t
\]
the associated characteristic polynomial is:
\[
z^p - a_1 z^{p-1} - a_2 z^{p-2} - \cdots - a_p = 0
\]

\textbf{Stationarity condition:} all $p$ roots of this equation must lie
\emph{inside} the complex unit circle, i.e., $|z_i| < 1$ for every root $z_i$.

\textit{Intuition:} this is a generalization of the AR(1) condition $|a_1| < 1$.
For AR(1), the characteristic equation is $z - a_1 = 0 \Rightarrow z = a_1$,
and $|z| < 1$ reduces immediately to $|a_1| < 1$.

\subsection{AR(1) as a Markov Process}
AR(1) has a special structural property: $X_t$ depends on history \emph{only}
through $X_{t-1}$. Given $X_{t-1}$, the earlier values $X_{t-2}, X_{t-3}, \ldots$
carry no additional information about $X_t$. This is the \textbf{Markov property}.

\begin{itemize}
    \item An AR(1) is a \textbf{first-order Markov process}: $P(X_t \mid X_{t-1},
      X_{t-2}, \ldots) = P(X_t \mid X_{t-1})$.
    \item More generally, an AR($p$) can always be rewritten as a $p$-dimensional
      AR(1) process (the \textit{companion form}) by stacking the last $p$ lags
      into a state vector.
\end{itemize}


%------------------------------------------------------------
\section{Yule-Walker Equations: Estimating AR Coefficients from ACF}
%------------------------------------------------------------

\subsection{Derivation}
Multiply both sides of the AR($p$) equation by $X_{t-k}$ and take expectations.
Using stationarity, the covariance $v(k) = E[(X_t - \mu)(X_{t-k} - \mu)]$ satisfies
a system of equations for $k = 1, 2, \ldots, p$:
\[
r(k) = a_1 r(k-1) + a_2 r(k-2) + \cdots + a_p r(k-p), \qquad k = 1, \ldots, p
\]
where $r(k) = v(k)/v(0)$ is the lag-$k$ autocorrelation.

\subsection{Matrix Form}
Writing all $p$ equations simultaneously:
\[
\begin{pmatrix} r(1) \\ r(2) \\ \vdots \\ r(p) \end{pmatrix}
=
\begin{pmatrix}
  r(0)   & r(1)   & \cdots & r(p-1) \\
  r(1)   & r(0)   & \cdots & r(p-2) \\
  \vdots & \vdots & \ddots & \vdots \\
  r(p-1) & r(p-2) & \cdots & r(0)
\end{pmatrix}
\begin{pmatrix} a_1 \\ a_2 \\ \vdots \\ a_p \end{pmatrix}
\]

This is a \textbf{linear system} --- given the sample ACF values $\hat{r}(0),
\hat{r}(1), \ldots, \hat{r}(p)$, the AR coefficients $a_1, \ldots, a_p$ can be
solved directly (Toeplitz matrix, efficient inversion algorithms exist).

\textit{Practical value:} You can estimate AR coefficients using only the
correlogram (ACF plot) without fitting the model by OLS. Both approaches
should give similar results for large samples.


%------------------------------------------------------------
\section{MA($q$) Processes}
%------------------------------------------------------------

\subsection{Definition}
A \textbf{Moving Average of order $q$}, MA($q$), expresses $X_t$ as a weighted
sum of \emph{current and past noise terms}:
\[
X_t = \mu + \varepsilon_t + b_1 \varepsilon_{t-1} + b_2 \varepsilon_{t-2}
      + \cdots + b_q \varepsilon_{t-q}
\]
where $\varepsilon_t \sim \text{WN}(0, \sigma^2)$ (white noise: i.i.d., zero mean,
constant variance).

\subsection{ACF of an MA($q$)}
The ACF of an MA($q$) is \textbf{finite}: it cuts off sharply to zero after
lag $q$:
\[
r(k) = 0 \quad \text{for all } k > q
\]
This is the key diagnostic: a sharp ACF cutoff identifies MA order directly,
whereas AR produces gradually decaying ACF.

\subsection{Invertibility}
An MA($q$) process is \textbf{invertible} if it can be written as an AR($\infty$)
process with summable coefficients. The invertibility condition mirrors the AR
stationarity condition: the roots of the MA characteristic equation
$1 + b_1 z + b_2 z^2 + \cdots + b_q z^q = 0$ must lie \emph{outside} the unit
circle.

\textit{Why it matters:} Non-invertible MA processes lead to non-unique
parameter estimates. Invertibility is the condition that ensures the model can
actually be identified and estimated from data.


%------------------------------------------------------------
\section{ARMA($p$,$q$) and ARIMA($p$,$d$,$q$) Models}
%------------------------------------------------------------

\subsection{ARMA($p$,$q$): Combining AR and MA}
An AR($p$) model uses only past $X$ values; an MA($q$) model uses only past
noise values. An \textbf{ARMA($p$,$q$)} uses both:
\[
X_t = a_1 X_{t-1} + \cdots + a_p X_{t-p}
    + \varepsilon_t + b_1 \varepsilon_{t-1} + \cdots + b_q \varepsilon_{t-q}
\]

\textbf{Backward operator notation.} Define the lag operator $B$ by
$B^j X_t = X_{t-j}$. Then:
\begin{align*}
A(B) &= 1 - a_1 B - a_2 B^2 - \cdots - a_p B^p  \quad\text{(AR polynomial)} \\
C(B) &= 1 + b_1 B + b_2 B^2 + \cdots + b_q B^q  \quad\text{(MA polynomial)}
\end{align*}
The ARMA model becomes the compact expression $A(B)X_t = C(B)\varepsilon_t$.

\subsection{ARIMA($p$,$d$,$q$): Handling Non-Stationarity}
Real-world series (prices, GDP, cumulative counts) often have a trending mean
--- they are \emph{not} stationary. The solution is \textbf{differencing}.

\begin{itemize}
    \item \textbf{First difference} ($d = 1$): replace $X_t$ with
      $\Delta X_t = X_t - X_{t-1}$. This removes a linear trend.
    \item \textbf{Second difference} ($d = 2$): difference again to remove
      quadratic drift. Rarely needed beyond $d = 2$ in practice.
\end{itemize}

An \textbf{ARIMA($p$,$d$,$q$)} model fits an ARMA($p$,$q$) to the
$d$-times-differenced series. Using the backward operator, differencing is
$(1 - B)^d$, giving:
\[
A(B)(1-B)^d X_t = C(B)\varepsilon_t
\]

\subsection{Wold's Decomposition Theorem}
Any zero-mean covariance-stationary process can be written as:
\[
X_t = \sum_{j=0}^{\infty} \psi_j \varepsilon_{t-j} + \kappa_t, \quad
\psi_0 = 1,\quad \sum_{j=0}^{\infty} \psi_j^2 < \infty
\]
where $\varepsilon_t$ is white noise and $\kappa_t$ is a perfectly predictable
deterministic component.

\textit{Practical implication:} Every stationary process is \emph{in principle}
an MA($\infty$). ARMA models are \textbf{parsimonious approximations} to this
infinite representation --- using a small number of AR and MA coefficients
to capture the same dynamics.


%------------------------------------------------------------
\section{Box-Jenkins Methodology: Systematic Model Selection}
%------------------------------------------------------------

The Box-Jenkins framework is the standard workflow for fitting ARIMA models.
It has three steps:

\begin{enumerate}
    \item \textbf{Transform to stationarity.}
    \begin{itemize}
        \item Check for trends and non-constant variance visually or via a
          unit root test (e.g., Augmented Dickey-Fuller).
        \item Difference until the series looks stationary ($d$ = number of
          differences).
        \item Apply log or square-root transformation to stabilize variance
          if needed.
    \end{itemize}

    \item \textbf{Estimate ARMA on the stationary series.}
    \begin{itemize}
        \item Inspect the ACF (identifies MA order $q$: sharp cutoff at lag $q$).
        \item Inspect the PACF (identifies AR order $p$: sharp cutoff at lag $p$).
        \item Fit candidate ARMA($p$, $q$) models; compare with AIC/BIC to
          balance fit against model complexity.
    \end{itemize}

    \item \textbf{Diagnose residuals (model checking).}
    \begin{itemize}
        \item Compute the ACF of the residuals.
        \item Apply the Ljung-Box test: residuals should be \textbf{white
          noise} (no remaining autocorrelation).
        \item A significant Ljung-Box $p$-value means the model is
          \emph{inadequate} --- return to step 2 and revise.
        \item An insignificant $p$-value (large $p$) means the model
          captures all the structure; residuals are unpredictable noise.
    \end{itemize}
\end{enumerate}


%------------------------------------------------------------
\section{White Noise Testing: The Ljung-Box Statistic}
%------------------------------------------------------------

\subsection{Motivation}
After fitting an AR/ARMA/ARIMA model, you test whether the residuals are
\emph{white noise}. If autocorrelation remains, the model missed something.

\subsection{The Statistic}
Let $\hat{r}_k$ be the sample autocorrelation of the residuals at lag $k$,
computed from a series of length $T$. The \textbf{Ljung-Box} statistic is:
\[
LB = T(T+2) \sum_{k=1}^{p} \frac{\hat{r}_k^2}{T - k}
\]

Under the null hypothesis that the residuals are white noise, $LB$ follows a
\textbf{chi-squared distribution} with $p$ degrees of freedom:
\[
LB \;\sim\; \chi^2(p)
\]

\begin{itemize}
    \item \textbf{Large $LB$} (small $p$-value): reject white noise hypothesis
      --- residuals still contain structure. The model needs improvement.
    \item \textbf{Small $LB$} (large $p$-value): fail to reject --- residuals
      look like white noise. The model is adequate.
\end{itemize}

The significance threshold $p = 0.05$ is the typical cutoff.


%------------------------------------------------------------
\section{Worked Example: Fitting AR(2) to AAA Bond Yield}
%------------------------------------------------------------

\subsection{Series and Identification}
The series is the \textbf{Moody's AAA corporate bond yield} (monthly data).
After inspecting the ACF and PACF, the analyst identifies an \textbf{AR(2)}
as the candidate model.

\subsection{Estimated Model}
After estimation, the fitted AR(2) is:
\[
X_t = c + 1.337606\, X_{t-1} - 0.383890\, X_{t-2} + \varepsilon_t
\]

\subsection{Stationarity Check via Characteristic Roots}
The characteristic equation $z^2 - 1.337606\,z + 0.383890 = 0$ has roots:
\[
z_1 \approx -0.417, \qquad z_2 \approx -0.920
\]
Both satisfy $|z_i| < 1$ (inside the unit circle), so the fitted model is
\textbf{stationary}.

\subsection{Residual Diagnostics}
Applying the Ljung-Box test to the residuals (with $p = 18$ lags):
\[
LB = 6.79, \qquad p\text{-value} = 96.3\%
\]
The $p$-value is very large --- we \textbf{fail to reject} the null of white
noise. Conclusion: the AR(2) residuals contain no detectable remaining structure.
The model is adequate.

\begin{center}
\begin{tabular}{@{}ll@{}}
\toprule
\textbf{Step} & \textbf{Result} \\
\midrule
Model         & AR(2): $X_t = c + 1.3376 X_{t-1} - 0.3839 X_{t-2} + \varepsilon_t$ \\
Char.\ roots  & $-0.417,\; -0.920$ (both $< 1$) $\Rightarrow$ stationary \\
Ljung-Box     & $LB = 6.79$, $p = 96.3\%$ $\Rightarrow$ residuals are white noise \\
Verdict       & Model is adequate; no further iteration needed \\
\bottomrule
\end{tabular}
\end{center}


%------------------------------------------------------------
\section{Missing Data: Strategies (Preview / HW Discussion)}
%------------------------------------------------------------

A student question during HW1 discussion prompted this in-class treatment.
When data has missing values (coded as $-999$ or blank), the right answer is always
\textbf{``it depends.''} The questions to ask first:

\begin{enumerate}
    \item \textbf{What fraction is missing?} --- 1\% vs.\ 50\% require completely
      different responses.
    \item \textbf{Which variables are missing?} --- High-importance features need
      more careful treatment than low-importance ones.
    \item \textbf{How are they missing?} --- Randomly distributed across the
      dataset (easier to handle) vs.\ concentrated in a specific range or
      correlated with another variable (more serious).
\end{enumerate}

\textbf{Strategy options} (in rough order of sophistication):

\begin{itemize}
    \item \textbf{Delete rows:} Appropriate only when missing fraction is very
      small (e.g., $<5\%$). Deleting a majority of data destroys information.
    \item \textbf{Mean imputation:} Replace missing values with the column average.
      Simple, better than nothing, but introduces bias toward the center.
    \item \textbf{XGBoost sparsity awareness:} Let XGBoost handle missing values
      natively by learning the optimal default direction. Converts missing
      indicators to empty rather than $-999$ (which XGBoost treats as a real
      value).
    \item \textbf{Conditional imputation (posterior estimate):} For each record,
      use all available features to estimate the missing value --- the most
      principled approach, but most complex.
    \item \textbf{Synthetic data augmentation:} When data is \emph{scarce} (not
      just missing), generate synthetic samples or apply data augmentation to
      increase effective training size.
\end{itemize}

\textit{For imbalanced data} (small signal, large background): Addressed next
week in detail. Options include resampling the minority class, adjusting class
weights, or using evaluation metrics (AUC, F1) that are robust to imbalance
rather than raw accuracy.

\end{document}
